% =====================================================================
%  1bat-ud6-6.tex  —  SESSIÓ 6: Gimnàstica de conjunts: probabilitats condicionades
%  (sense preàmbul: s'inclou des de main.tex)
% =====================================================================
\horatitol{Sessió 6 — Gimnàstica de conjunts: probabilitats condicionades}%
{Mecanitzem la fórmula de la condicionada: calcular-la en tots dos sentits i comprovar si dos successos són independents.}

\subsection{Idea clau}

Completem l'Activitat 3 amb la \textbf{mecanització} de les probabilitats
condicionades, ja introduïdes de manera intuïtiva a la sessió 5. Avui, amb exercicis
breus i simbòlics, automatitzem la fórmula i una comprovació clau: si dos successos són
\textbf{independents}.

\begin{idea}
\textbf{Condicionada en tots dos sentits i independència}
\begin{itemize}[leftmargin=1.4em, itemsep=3pt]
  \item \textbf{La fórmula, en els dos sentits:}
        \[
        P(A\mid B)=\frac{P(A\cap B)}{P(B)},\qquad
        P(B\mid A)=\frac{P(A\cap B)}{P(A)}.
        \]
        \textbf{Compte:} no és el mateix $P(A\mid B)$ que $P(B\mid A)$; canvia el
        denominador.
  \item \textbf{Test d'independència:} $A$ i $B$ són independents si
        \[
        P(A\cap B)=P(A)\per P(B).
        \]
        Si es compleix, saber-ne un no canvia la probabilitat de l'altre.
\end{itemize}
\end{idea}

\subsection{Exemples}

\begin{exemple}[la mateixa intersecció, dues condicionades]
Se sap que $P(A)=0,6$, $P(B)=0,5$ i $P(A\cap B)=0,2$. Les dues condicionades:
\[
P(A\mid B)=\frac{0,2}{0,5}=0,4,\qquad
P(B\mid A)=\frac{0,2}{0,6}\approx 0,33.
\]
Surten \textbf{diferents} perquè el denominador (la condició) és diferent. La pregunta
«$A$ sabent $B$» no és la mateixa que «$B$ sabent $A$».
\end{exemple}

\begin{exemple}[comprovar la independència]
Amb $P(A)=0,5$, $P(B)=0,4$ i $P(A\cap B)=0,2$, són independents? Apliquem el test:
\[
P(A)\per P(B)=0,5\per 0,4=0,2=P(A\cap B).
\]
Com que coincideixen, $A$ i $B$ \textbf{són independents}: saber-ne un no canvia la
probabilitat de l'altre (de fet, $P(A\mid B)=\frac{0,2}{0,4}=0,5=P(A)$).
\end{exemple}

\subsection{Exercicis resolts}

\begin{exresolt}[les dues condicionades i el test]
Se sap que $P(A)=0,4$, $P(B)=0,5$ i $P(A\cap B)=0,1$. Calcula $P(A\mid B)$ i
$P(B\mid A)$, i digues si $A$ i $B$ són independents.
\solucio
\[
P(A\mid B)=\frac{0,1}{0,5}=0,2,\qquad P(B\mid A)=\frac{0,1}{0,4}=0,25.
\]
\textbf{Test d'independència:} $P(A)\per P(B)=0,4\per 0,5=0,2$, però $P(A\cap B)=0,1$.
Com que $0,2\neq 0,1$, \textbf{no} són independents (són \textbf{dependents}).
\resposta{$P(A\mid B)=0,2$; $P(B\mid A)=0,25$; són dependents.}
\end{exresolt}

\begin{exresolt}[de la condicionada a la intersecció]
Se sap que $P(B)=0,4$ i $P(A\mid B)=0,3$. Calcula $P(A\cap B)$.
\solucio
Aïllem la intersecció de la fórmula de la condicionada
($P(A\mid B)=\frac{P(A\cap B)}{P(B)}$), multiplicant per $P(B)$:
\[
P(A\cap B)=P(A\mid B)\per P(B)=0,3\per 0,4=0,12.
\]
\resposta{$P(A\cap B)=0,12$.}
\end{exresolt}

\subsection{Exercicis per fer}

\begin{exercicis}
\begin{exlist}
  \item Amb $P(A)=0,5$, $P(B)=0,6$ i $P(A\cap B)=0,3$, calcula $P(A\mid B)$ i
        $P(B\mid A)$.
  \item Amb $P(A)=0,2$, $P(B)=0,5$ i $P(A\cap B)=0,1$, comprova si $A$ i $B$ són
        independents (test $P(A)\per P(B)$).
  \item Amb $P(B)=0,5$ i $P(A\mid B)=0,4$, calcula $P(A\cap B)$.
  \item Amb $P(A)=0,7$ i $P(B\mid A)=0,2$, calcula $P(A\cap B)$.
  \item \textbf{(Compte amb l'ordre)} Se sap que $P(A\cap B)=0,15$, $P(A)=0,3$ i
        $P(B)=0,5$. Calcula $P(A\mid B)$ i $P(B\mid A)$ i comprova que surten diferents.
        Quin és més gran i per què?
  \item \textbf{(Raonament)} Si dos successos són \textbf{independents}, què val
        $P(A\mid B)$ en funció de $P(A)$? Comprova-ho amb l'exemple
        $P(A)=0,5$, $P(B)=0,4$, $P(A\cap B)=0,2$.
\end{exlist}
\end{exercicis}

% ---------------------------------------------------------------------
\fullsolucions{Sessió 6}
\begin{solucions}
\begin{sollist}
  \item $P(A\mid B)=\dfrac{0,3}{0,6}=0,5$; $P(B\mid A)=\dfrac{0,3}{0,5}=0,6$.
  \item $P(A)\per P(B)=0,2\per 0,5=0,1=P(A\cap B)$: coincideixen, de manera que $A$ i
        $B$ \textbf{són independents}.
  \item $P(A\cap B)=P(A\mid B)\per P(B)=0,4\per 0,5=0,2$.
  \item $P(A\cap B)=P(B\mid A)\per P(A)=0,2\per 0,7=0,14$.
  \item $P(A\mid B)=\dfrac{0,15}{0,5}=0,3$; $P(B\mid A)=\dfrac{0,15}{0,3}=0,5$. Surten
        diferents. És més gran $P(B\mid A)=0,5$ perquè té un denominador més petit
        ($P(A)=0,3<P(B)=0,5$): dividir per un nombre més petit dóna un resultat més gran.
  \item Si són independents, $P(A\mid B)=P(A)$ (saber $B$ no canvia res). Comprovació:
        $P(A\mid B)=\dfrac{0,2}{0,4}=0,5=P(A)$. Es compleix.
\end{sollist}
\end{solucions}
